% Post-checkpoint08 source. No integrated file is modified.
\subsection{The normalized configuration source of the translation comparison}

\begin{result}{A.4.46}{Proved}{The actual translation action has a unique fixed lift and trivial configuration-kernel centralizer}
Let $T=\mathbb R^2/\mathbb Z^2$ and index nine ordered
points by $H=\mathbb F_3^2$. The simultaneous-translation
quotient of their configuration space is the actual space
\[
 X=\operatorname{Conf}_9(T)/T
  \xrightarrow{\sim}\operatorname{Conf}_8(T\setminus\{0\}),
 \qquad(z_h)\longmapsto(z_h-z_0)_{h\ne0}.
                                                               \tag{NCF.1}
\]
Choose the grid configuration $z_h=h/3$ as basepoint,
and let $K_0=\pi_1(X)$.
The finite translation-label action is exactly
\[
                    (a\cdot z)_h=z_{h+a}-z_a.
                                                               \tag{NCF.2}
\]
Its fixed locus has exactly $48$ points. The action
lifts uniquely to the universal cover while fixing
the chosen grid lift; that lifted action has precisely
one fixed point. Consequently
\[
                         C_{K_0\rtimes H}(H)\cap K_0=1.
                                                               \tag{NCF.3}
\]
Here $K_0$ is the actual pure configuration kernel
in the nine-point torus mapping class group, and
the semidirect action is its conjugation by the
named translation subgroup. This result is an
ordinary source statement; it is not a claim that
(NCF.3) holds in every finite quotient.
\end{result}

\begin{proof}
The homeomorphism
$\operatorname{Conf}_9(T)\cong T\times X$
and its inverse are respectively
$(z_h)\mapsto(z_0,(z_h-z_0)_{h\ne0})$ and
$(t,(w_h))\mapsto(t+w_h)$, with $w_0=0$.
Relabelling by $h\mapsto h+a$ and returning the
new zero point to the origin gives (NCF.2).
Direct substitution proves the action law for every
$a,b\in H$. Thus it is the actual label-translation
action, with no forgotten basepoint adjustment.
For its precise group parameter, write $T_t(z)=z+t$.
Conjugating a point-pushing path by $T_{-a/3}$ gives
the normalized path $z_{h+a}-z_a$, whereas conjugation
by $T_{a/3}$ gives $z_{h-a}-z_{-a}$.
Thus in (NCF.2) the element $a$ of the abstract label
group is identified with the actual translation
$T_{-a/3}$.  The comparison with a convention using
$T_{a/3}$ is the explicit group automorphism
$a\mapsto-a$.  This is the same embedded translation
subgroup, and this scalar label automorphism commutes
with every linear normalizer action.
The evaluation fibration of the identity component
of the torus diffeomorphism group gives the same
configuration kernel: its translation factor is
the displayed $T$ and its quotient is (NCF.1).
This identifies the action on $K_0$ with ordinary
conjugation by the marked translations.

A configuration is fixed by (NCF.2) precisely when
\[
                 z_{h+a}=z_h+z_a\quad(h,a\in H).
                                                               \tag{NCF.4}
\]
Thus $h\mapsto z_h$ is a homomorphism $H\to T$.
The distinct-point condition makes it injective,
and its image is necessarily all of $T[3]$.
Conversely every isomorphism $H\to T[3]$ gives a
fixed configuration. There are
$(3^2-1)(3^2-3)=48$ such isomorphisms.
In particular this fixed locus is discrete.

The space $X$ is a smooth real manifold of dimension
$16$. Successively forgetting the last point in
$\operatorname{Conf}_j(T\setminus\{0\})$ has fibre
the torus with $j$ points removed, of free
fundamental rank $j+1$, for $1\leq j\leq8$.
Local disjoint discs around the retained points
give the local trivializations of these actual
Fadell--Neuwirth fibrations. The fibres and the
initial punctured torus are aspherical. Their
homotopy exact sequences inductively show that
every higher homotopy group of $X$ vanishes.
The universal cover is a simply connected CW-type
manifold with all homotopy groups zero, hence is
contractible.

Each element of $H$ has a unique lift fixing the
selected point of the universal cover. Products
fix that point and project to the correct products,
so these lifts form an actual action of $H$.
The action and its lifts are smooth, hence locally
linear. Smith theory applies to this finite
$3$-group acting on the contractible, finite-dimensional
manifold: its fixed set is nonempty and
modulo-three acyclic, in particular connected.
An explicit opened reference with these hypotheses
is Edmonds \cite[printed pp.~273--274, Proposition~3.2.1,
Theorem~3.3.1 and Proposition~3.3.2]{ET_Edmonds2012}.

The fixed set upstairs projects into the $48$-point
set (NCF.4). Since the covering map is a local
homeomorphism, this upstairs fixed set is discrete.
Its connectedness therefore makes it exactly the
one selected lift. A deck transformation commuting
with $H$ carries that fixed point to another fixed
point, hence fixes it. A deck transformation fixing
a point is the identity. This proves (NCF.3).
\end{proof}
